\section{Physics Simulation: the Gazebo Integration}
\label{sec:gazebo}

The research pipeline described so far lives entirely in an idealized
simulation harness: positions are coordinates in an array, links appear and
disappear according to an analytic radio model, and outcomes are sampled
from a hidden reliability field. That harness is the right tool for the
statistical experiments of the previous chapters, because it is fast,
seeded, and exactly reproducible. It is not, however, a physics engine,
and it does not exercise the integration questions that any real
deployment must answer: how the pipeline receives live odometry from a
running simulator, how it is packaged so that a reviewer can reproduce it
with one command, and how its decisions can be observed while a world is
being simulated. This chapter describes how we implemented exactly that
in Gazebo: a containerized integration in which four differential-drive
rovers are simulated by Gazebo Harmonic, their odometry is bridged into
ROS~2 Jazzy, and the full research pipeline (QUBO pruning, trust-aware
GNN scoring, task assignment, and trust updates) runs inside a ROS~2 node
against that live stream, publishing its decisions as JSON status
messages once per second.

We emphasize at the outset that this chapter reports an integration, not
a new experimental result. The physics simulator currently contributes
geometry and ideal odometry; the cognitive pipeline on top is the same
one audited in the previous chapters, with the same delivery model and
the same trust dynamics. Section~\ref{sec:gazebo:limitations} states the
resulting limitations plainly, together with the remedy we plan for
each.

\subsection{Software Stack and Deployment}
\label{sec:gazebo:stack}

The integration is built on Gazebo Harmonic (the \code{gz-sim}
generation of the simulator), using SDF~1.9 as the world and model
description format, paired with ROS~2 Jazzy as the middleware. The two
systems are deliberately kept at arm's length: Gazebo simulates the world
and exposes state over its own transport, and a narrow bridge forwards
exactly the topics the planner needs into ROS~2. This separation mirrors
how the pipeline would sit on a real fleet, where the estimator and the
planner are different processes with a well-defined interface between
them.

Everything runs in Docker. The image is based on the ROS~2 Jazzy
distribution, with the Gazebo ROS bridge, the navigation and geometry
message packages, and NumPy installed on top. The whole project is copied
into the image so that the research pipeline is importable as plain
Python modules; we deliberately do not wrap the research code as ROS
packages, because the pipeline is middleware-agnostic by design and the
ROS node is only a thin adapter around it. A compose file defines a
single service that builds this image and starts both the simulator and
the planner. A launch script guards the build: it refuses to proceed with
less than 8~GB of free disk, since the base image plus the simulator
toolchain is large, and a truncated build produces confusing failures far
from the real cause. After the check, the script builds and runs the
container.

The Gazebo server runs headless inside the container. A \code{-{}-gui}
flag exists for local debugging, but it requires display forwarding from
the host and is not part of the default path. No ports are published:
the demo is self-contained, and all observation happens through the log
stream and the two status topics described below.

\subsection{Interfaces: Bridging Gazebo and the Planner}
\label{sec:gazebo:interfaces}

A single parameter bridge connects Gazebo transport to ROS~2. It
forwards exactly four topics, one per rover, and nothing else:
\code{/model/rover\_N/odometry} for $N$ in $0$ to $3$, each carrying
\code{nav\_msgs/Odometry}. We chose to bridge only odometry, rather than
the full simulator state, for two reasons. First, odometry is the
minimal state the planner actually consumes: positions and velocities of
the fleet. Second, a narrow bridge keeps the data contract explicit and
auditable; anything the planner knows about the world arrives on one of
these four streams.

In the opposite direction, the planner node publishes two JSON status
topics, both of type \code{std\_msgs/String}: \code{/quantum\_gnn/status}
carries a per-cycle summary of the whole pipeline (cycle index, edge
counts, solver time, assignment, and trust vector), and
\code{/quantum\_gnn/assignments} carries the current task assignment as
structured JSON. Table~\ref{tab:gazebo:interfaces} lists the complete
interface.

\begin{table}[t]
  \centering
  \caption{ROS~2 and Gazebo interfaces of the integration. Direction is
    given relative to the planner node. The four odometry topics are
    bridged from Gazebo transport; the two status topics are native
    ROS~2 publications.}
  \label{tab:gazebo:interfaces}
  \begin{tabular}{@{}llll@{}}
    \toprule
    Topic & Message type & Direction & Rate \\
    \midrule
    \code{/model/rover\_0/odometry} & \code{nav\_msgs/Odometry} & subscribe (bridged) & 20\,Hz \\
    \code{/model/rover\_1/odometry} & \code{nav\_msgs/Odometry} & subscribe (bridged) & 20\,Hz \\
    \code{/model/rover\_2/odometry} & \code{nav\_msgs/Odometry} & subscribe (bridged) & 20\,Hz \\
    \code{/model/rover\_3/odometry} & \code{nav\_msgs/Odometry} & subscribe (bridged) & 20\,Hz \\
    \code{/quantum\_gnn/status} & \code{std\_msgs/String} (JSON) & publish & 1\,Hz \\
    \code{/quantum\_gnn/assignments} & \code{std\_msgs/String} (JSON) & publish & 1\,Hz \\
    \bottomrule
  \end{tabular}
\end{table}

\begin{figure}[t]
  \centering
  \resizebox{0.95\textwidth}{!}{%
  \begin{tikzpicture}[
      font=\small,
      box/.style={draw, rounded corners=2pt, align=center,
                  minimum height=9mm, inner sep=5pt},
      stage/.style={draw, align=center, minimum height=7mm,
                    inner sep=4pt, fill=black!4},
      flow/.style={-{Stealth[length=2.5mm]}, thick},
      lbl/.style={font=\footnotesize, midway, above, align=center},
    ]
    \node[box] (gz) {Gazebo Harmonic\\ODE physics, 1000\,Hz};
    \node[box, right=16mm of gz] (bridge) {parameter bridge\\Gazebo transport $\to$ ROS~2};
    \node[box, right=16mm of bridge] (planner) {planner node\\ROS~2 Jazzy, 1\,Hz timer};
    \node[box, below=14mm of planner] (topics) {\code{/quantum\_gnn/status}\\\code{/quantum\_gnn/assignments}};

    \node[stage, above=10mm of planner, xshift=-30mm] (s1) {1.\ candidate graph\\from odometry};
    \node[stage, right=4mm of s1] (s2) {2.\ QUBO pruning\\projected annealer};
    \node[stage, right=4mm of s2] (s3) {3.\ GNN scoring\\$\times$ widest-path factor};
    \node[stage, right=4mm of s3] (s4) {4.\ assignment, outcomes,\\trust update};

    \begin{scope}[on background layer]
      \node[draw, dashed, inner sep=6pt,
            fit=(s1)(s2)(s3)(s4), label={[font=\footnotesize]above:pipeline stages (one cycle)}] (stages) {};
    \end{scope}

    \draw[flow] (gz) -- node[lbl]{\code{/model/rover\_N/odometry}\\$N=0,\dots,3$, 20\,Hz} (bridge);
    \draw[flow] (bridge) -- node[lbl]{\code{nav\_msgs/Odometry}} (planner);
    \draw[flow] (planner) -- node[lbl, right]{JSON, 1\,Hz} (topics);
    \draw[flow, dashed] (planner.north) -- (stages.south -| planner.north);
    \draw[flow, dashed] (stages.south -| s2.south) -- ++(0,-4mm) -| (planner.north);
  \end{tikzpicture}}
  \caption{Dataflow of the Gazebo integration. Odometry flows from the
    physics server through the parameter bridge into the planner node,
    which executes the six-stage pipeline cycle once per second and
    publishes its decisions as JSON status topics.}
  \label{fig:gazebo:dataflow}
\end{figure}

\fillme{screenshot of the Gazebo world with the four rovers}

\subsection{The Planner Node Cycle}
\label{sec:gazebo:cycle}

The planner node runs one pipeline cycle per second on a wall-clock
timer. The number of cycles defaults to twelve and is set by an
environment variable, so a short demonstration run finishes on its own
and the container exits. Each cycle executes the same six stages as the
research harness, fed by live odometry instead of array positions.

\paragraph{1. Candidate graph and edge features.}
From the latest odometry of the four rovers we build the candidate graph:
the complete graph on the fleet, restricted to edges whose endpoints are
within the communication range of $5.0$\,m. Each candidate edge carries
the same seven features as in the research pipeline, computed
deterministically from the current state. Trust is the minimum of the
two endpoint trust values, so a link is only as trustworthy as its
weaker endpoint. Proximity is $1 - d/r$ clipped to $[0,1]$, where $d$ is
the inter-rover distance and $r$ the communication range. Link quality
is $0.85$ times proximity plus $0.15$: an analytic radio model, chosen
deliberately so that container runs are reproducible and comparable
across machines, with a floor of $0.15$ so that an in-range link is
never hopeless. Processing cost is $0.10 + 0.50 \cdot (1 -
\text{proximity})$, so relaying across a long, weak link is priced
higher than across a short one. Battery is the minimum of the two
endpoint battery levels, freshness is $1$ (odometry is live), and task
criticality is $0.5$ for all edges.

\paragraph{2. Constrained QUBO pruning.}
The candidate graph is pruned by the constrained QUBO formulation of the
previous chapters: a degree cap of three per node, a global processing
budget of $25$ units, all four rovers required to remain connected, and
rover~0 designated as the root toward which connectivity is enforced by
the single-commodity flow constraints. The QUBO is solved by the
production projected annealer, configured with a fast demonstration
schedule: initial temperature $0.5$, final temperature $0.05$, cooling
rate $0.80$, two steps per temperature level, and two restarts. The kept
edge set of the previous cycle is supplied as a warm start, so in steady
state the annealer refines a good solution rather than searching from
scratch; recall that the restarts parameter counts the warm start, which
is why the minimum usable value here is two. Solve time is measured and
reported every cycle. Because the annealer is heuristic, the planner
treats infeasibility as a first-class event: if no feasible pruning is
found, it falls back to keeping all candidate edges, and the fallback is
logged and flagged in the status topic rather than silently absorbed.

\paragraph{3. GNN scoring with path factors.}
The trust-aware GNN scores the four rovers on the \emph{kept} graph
only, exactly as in the research harness: scores computed on pruned-away
structure would credit rovers for connectivity they cannot actually use.
Each raw score is then multiplied by a widest-path factor from the root:
the maximum, over all paths from rover~0 to the rover in the kept graph,
of the minimum link quality along the path. This is the same
bottleneck-delivery model used throughout the project, so a rover that
is well scored but reachable only through a weak link is discounted
accordingly.

\paragraph{4. Task assignment.}
Two tasks are assigned per cycle, with criticalities $0.7$ and $0.4$,
by the capacity-aware greedy matcher: each rover executes at most one
task per cycle, and the higher-criticality task is placed first.

\paragraph{5. Outcomes and trust updates.}
Task outcomes are sampled from hidden per-rover reliabilities of $0.92$,
$0.82$, $0.75$, and $0.66$, with delivery success drawn as a Bernoulli
trial in the rover's path factor, so mission success causally depends on
the topology the pruner kept, not merely on which rover was chosen.
Trust is updated only on execution outcomes, using an exponential moving
average with $\alpha = 0.55$ and a decay of $0.08$; rovers that received
no task produce no evidence, and their trust relaxes toward the prior
instead of being rewarded or punished for idleness. Computational loads
and battery levels evolve according to the same scripted dynamics as the
harness.

\paragraph{6. Status publication.}
Finally, the node publishes its JSON status: the cycle index, the number
of odometry messages received, the candidate and kept edge counts, the
pruner mode (normal or fallback), the measured solve time in
milliseconds, the assignment, and the current trust vector. This message
is the primary observability surface of the demo; a reviewer can watch
the pruner react to the fleet in real time by echoing one topic.

\fillme{screenshot or log excerpt of a live status message}

\subsection{World and Rover Model}
\label{sec:gazebo:world}

The world description defines a flat, empty arena of $30 \times 30$\,m,
lit by a single directional sun light with shadows enabled. Physics is
integrated by ODE at $1$\,ms steps, that is $1000$\,Hz, with a real-time
factor of $1.0$. The arena is deliberately featureless: the integration
is a plumbing test for the pipeline, not a navigation benchmark, and
obstacles would add a sensing problem we have intentionally not
introduced yet.

Four rovers are spawned at $(-3, 0)$, $(-1, 0.4)$, $(1, 0.1)$, and
$(3, 0.5)$ with small yaw offsets, all well within mutual communication
range, so the initial candidate graph is complete. The rover model is a
minimal differential-drive platform: a box chassis of $0.9 \times 0.65
\times 0.25$\,m with mass $3.0$\,kg, two cylindrical wheels of radius
$0.16$\,m and mass $0.25$\,kg each, placed at $\pm 0.38$\,m laterally
and attached by revolute joints. The Gazebo differential-drive plugin
drives the wheel joints with a wheel separation of $0.76$\,m and a wheel
radius of $0.16$\,m, and publishes odometry at $20$\,Hz, which is the
stream the bridge forwards. The model carries no IMU, no lidar, and no
camera, and its odometry is ideal and noise-free. This is a conscious
choice for a first integration: every quantity the planner consumes is
either exactly known or explicitly scripted, so any unexpected behavior
in the pipeline can be attributed to the pipeline itself rather than to
sensor noise.

\subsection{Limitations and Planned Remedies}
\label{sec:gazebo:limitations}

In keeping with the audit stance of this project, we state what the
integration does \emph{not} yet do, and what we plan to do about it.

First, \textbf{nothing publishes velocity commands}. The planner assigns
tasks, but no component translates an assignment into motion, so the
rovers hold their spawn poses for the entire run and Gazebo currently
contributes static geometry and an ideal odometry stream rather than
genuinely dynamic behavior. The positions the pipeline reasons over are
therefore constant across cycles, and only the scripted trust, load, and
battery dynamics change. The planned remedy is a thin executive layer
that maps each assignment to a target pose and publishes velocity
commands on the differential-drive plugin's command topic, closing the
loop between allocation and motion.

Second, \textbf{battery, load, capability, and link quality are
scripted, not sensed}. The radio model is analytic by design for
reproducibility, and the energy and load dynamics are the same
deterministic scripts as in the research harness. A real deployment
would estimate link quality from measured packet loss and round-trip
times and read battery state from the platform; we plan to replace the
analytic model with an estimator fed by the simulator's own topic
statistics as a first step in that direction, keeping the analytic model
as a reproducibility fallback.

Third, \textbf{the rover count is fixed at four in three places}: the
bridge's topic list, the planner's subscriptions, and the world
description all hard-code the fleet size, so adding a rover requires
three coordinated edits and silently drops the newcomer if one is
missed. The remedy is a single fleet-size parameter expanded into the
bridge configuration and the world description at launch time, with the
planner discovering its roster from the bridged topics.

Fourth, \textbf{there is no sim-time synchronization}: the planner ticks
on a one-second wall-clock timer, so if the simulator's real-time factor
drops below $1.0$ under load, pipeline cycles no longer align with
simulated seconds. The remedy is to drive the cycle from the simulated
clock topic instead of the wall clock, which the current node structure
accommodates with a one-line change to the timer source.

\subsection{Running the Demonstration}
\label{sec:gazebo:run}

The entire demonstration builds and runs with a single compose command
through the launch script, which performs the disk-space check, builds
the image, starts the Gazebo server and the bridge, and launches the
planner node. The \code{STEPS} environment variable sets the number of
pipeline cycles (twelve by default), after which the node reports its
final status and the container stops. No ports are published and no host
setup beyond Docker is required; the only artifacts are the container
log and the two JSON status topics, which together record every decision
the pipeline made and the state it made them from.
